Automated security assessments require an analyst to translate a natural-language
request into tool commands, run those commands, interpret their output, and
recover when a tool fails. Each step is manual, error-prone, and difficult to
automate end-to-end because an error message in the middle of a scan normally
stops the whole operation. This dissertation presents NeuroShell, an autonomous,
natural-language-driven orchestration framework for offensive security
operations that addresses this gap with four cooperating components.

The Intent Recognition Engine (Component 1) converts free-form security requests
into validated, machine-readable intent contracts. A small, offline, CPU-only
language model acts as a guarded natural-language gateway: a multi-stage
pipeline and deterministic validation stack reject malformed, out-of-scope and
adversarial outputs before any downstream action. The Dynamic Planner and
Retrieval-Augmented Generation Engine (Component 2) translates a structured
intent into a tool command for Nmap, Nikto or Gobuster, using retrieval over a
documentation knowledge base to ground command synthesis, and applies explicit
validation and safety checks. The Adaptive Execution and Error-Recovery Engine
(Component 3) runs every command in a locked-down Docker sandbox, and when a
command fails it labels the error, plans a fix with a local AI model, and retries
automatically, while refusing any command that fails its safety checks. The
AI-Driven Vulnerability Analysis Engine (Component 4) parses the execution
records produced by the security tools, enriches findings with CISA KEV, EPSS
and NVD intelligence, and prioritises vulnerabilities with an XGBoost-based
classifier explained by SHAP values.

NeuroShell keeps model inference entirely on the local machine, so no command,
error text or target ever leaves the deployment boundary. Evaluations measured
command-intent accuracy at 87.54\% for the planner, error-labelling macro-F1 of
99.1\% (and above 96\% under heavy message noise) for the recovery engine, and a
curated vulnerability-classification pipeline with strong discriminative
performance on a 155,852-record data set. Findings confirm that deterministic
validation, documentation grounding, sandboxed recovery and explainable
vulnerability scoring each contribute to a safer automated workflow. The
framework narrows the gap left by current AI copilots, which stop at tool
selection, execute without an autonomous error-recovery loop, or rely on
cloud-hosted models that conflict with strict data-governance requirements.

\sliitkeywords{Offensive Security, Natural Language Orchestration, Intent
Recognition, Retrieval-Augmented Generation, Error Recovery, Vulnerability
Analysis, Offline Language Models}